\documentclass[10pt]{article}
\usepackage[margin=20mm]{geometry}
\usepackage{amsmath,amssymb,mathtools,fontspec,unicode-math,xurl,hyperref}
\IfFontExistsTF{DejaVu Serif}{\setmainfont{DejaVu Serif}}{\setmainfont{Latin Modern Roman}}
\setmathfont{latinmodern-math.otf}
\setlength{\parindent}{0pt}\setlength{\parskip}{6pt}\sloppy
\title{Tautological cycles, local equations and nonzero certificates:
a version-specific mathematical review}
\author{AI-assisted technical review}
\date{}
\begin{document}
\maketitle
Abstract: The assigned publisher version is examined at the level of definitions, intersection multiplicities, Artin local equations, excess classes, Hodge integrals and the supplied computational evidence. The saved genus-six rank test leaves the essential nonzero assertion uncertified. Two demonstrated defects concern the incompatible Noether–Lefschetz definitions and a nonproper gluing in the core local-equation proof. The gluing admits a verified repair. These findings do not establish that the corrected main non-tautology theorem is false. Local arithmetic, sign and citation corrections, positive checks, and precise dependency and computation limits are also recorded.\par
Review status. The bounded proof and dependency checks described below are complete. Scientific decision is pending; human reviews are empty. The accompanying verification files record the independent focused checks and their completion status.\par
\section*{Source identity, hypotheses and proof structure}
The reviewed source is S. Canning, D. Oprea and R. Pandharipande, Tautological and non-tautological cycles on the moduli space of Abelian varieties, Inventiones mathematicae 242 (2025), 659–723, DOI https://doi.org/10.1007/s00222-025-01367-4 [main]. It is the assigned 65-page publisher PDF, with SHA256\par
`606c876e6cec594d5e44876c138d3f3855`
`49ca5e9254b5a8d10eafca15915dc3`.\par
Locators below are printed pages; PDF page is printed page minus \(658\). All 3,545 extracted lines, including references, were read to zero remaining. Relevant source and primary-dependency pages were also inspected as images. All fourteen files listed in the preserved source manifest were hashed again and matched. The current arXiv v3, dated 27 August 2025, was identified but not substituted for this PDF. The publisher records online publication on 9 September 2025. A bounded correction-heading lookup is not an exhaustive erratum search.\par
All Chow and cohomology groups are rational. The ambient \(\mathcal A_g\) is the smooth stack of principally polarized abelian varieties. The Torelli source \(\mathcal M_g^\mathrm{ct}\) is used for \(g\ge2\); the \(g=1\) product convention \(\mathcal A_0=\mathrm{point}\) is separate. The class \([\mathcal A_1\times\mathcal A_{g-1}]\) is the pushforward of the ordered product morphism, not automatically the reduced image with multiplicity one. This distinction matters in genus two.\par
Write \[\Delta_g=[\mathcal A_1\times\mathcal A_{g-1}]
       -\frac{g}{6|B_{2g}|}\lambda_{g-1},
 \qquad \delta_g=\operatorname{Tor}^*\Delta_g.\] The proof first identifies the only possible tautological expression for the product class, then proves that \(\delta_g\) lies in the radical of the \(\lambda_g\)-pairing. Local equations and excess-intersection calculations express it as a tautological class on \(\mathcal M_g^\mathrm{ct}\). Finally, the complete genus-six tautological-ring presentation and a computation are used to assert \(\delta_6\ne0\). The last step is essential: belonging to a one-dimensional radical does not by itself mean that a class generates that radical.\par
The genus-seven argument and the proposed general Noether–Lefschetz ring are distinguished from this main chain. The assertions explicitly called Conjectures or Speculation in the paper remain conjectural here. Conditional genus-eight and genus-nine computations are not promoted to unconditional results.\par
\section*{The essential genus-six nonzero test}
Problem 1 (Needs Evidence: the supplied test does not certify nonvanishing).  The main non-tautology conclusion and Theorem 5 require \(\delta_6\ne0\). The preserved `Delta6.ipynb` ends with a test that also succeeds when \(\delta_6=0\). An explicit nonzero certificate in the stated quotient is needed.\par
\subsection*{Exact evidence and its logical content}
The repository input is pinned to commit\par
`0a62f2dc9688fb0480967657ec9114321558085b`.\par
All seven notebooks were statically parsed and read, including all saved outputs: 98 cells in total, of which 88 are code, 9 markdown and 1 raw; 34 cells have saved outputs. No author code was executed. The relevant `Delta6.ipynb` has SHA256\par
`4b12741c197102d59b7999fb63d4c739f`
`559d4f47ee7a8a6e32fdd35bd0e5cf4`.\par
Cells 28, 32 and 33 contain, in the notebook’s notation,\par
{\ttfamily\footnotesize v = M.kernel().basis()[0]\\
Delta6 = list(-2730/691*lambda5 + ... )\\
A = matrix([Delta6,v])\\
A.rank()}\par
The final saved output is \(1\). The omitted expression in this display is the supplied sum of 24 named tree contributions; its full text is preserved in the static evidence extract, not omitted from the actual check.\par
Even granting the correct quotient basis, the correct matrix \(M\), and the nonzero kernel-basis vector \(v\), this proves only \[\delta_6\in\mathbb Qv.\] It does not prove that the scalar is nonzero. The exact countermodel \[\delta=(0,0),\qquad v=(1,0)\] has the same stacked-row rank \(1\). No independently displayed nonzero coordinate, nonzero scalar, single-vector rank or inequality \(\delta_6\ne0\) was found elsewhere in the seven supplied notebooks.\par
This is a logical insufficiency of the retained test. The notebook describes it as a proportionality check, whereas the paper separately asserts nonvanishing. The retained output does not establish that the authors inferred nonvanishing solely from rank one. The countermodel concerns that inference, not the actual assembled class. The nonvanishing claim therefore remains Needs Evidence in this bounded check; neither a zero class nor an incorrect saved rank has been demonstrated. The absence of author code would not itself be a defect; here there is an actual supplied test whose success does not distinguish the claimed conclusion from its negation.\par
\subsection*{The independent ring theorem does not close this step}
The cited Canning–Larson–Schmitt author PDF was obtained from the institutional URL in the bibliography [cls]. Its Theorem 7 and table give the complete presentation and the relevant dimensions, including the one-dimensional genus-six radical. Those are stronger than a numerical conjecture about relations. Nevertheless, completeness of a quotient and dimension one of its radical do not show that this particular assembled vector is nonzero.\par
The same dependency’s application paragraph identifies this specific class by referring back to the reviewed paper. Using that paragraph as an independent certificate would be circular. The actual cited author PDF is distinguished in the evidence ledger from the later 2026 arXiv version, which was not silently substituted.\par
A sufficient repair is small in output size but mathematically decisive: give a coordinate \(q\ne0\) of \(\delta_6\) in a certified quotient basis, exhibit \(\delta_6=qv\) with \(q\ne0\), or provide an equivalent exact nonzero witness whose relation reduction can be checked. The basis and relation version must be bound to that witness. The main conclusion remains needs evidence at this step; the report does not claim a counterexample to it.\par
\section*{The Noether–Lefschetz locus is not consistently specified}
Problem 2 (Major: exact rank, closure and the classification differ).  On p. 668 the locus \(\mathrm{NL}_g^r\) is defined by \(\operatorname{NS}(X)\cong\mathbb Z^r\), and its Zariski closure is then used. Theorem 13, pp. 669–670, attributes to Debarre–Laszlo a classification of the resulting rank-two locus. The primary theorem classifies the locus of rank greater than one, and its listed components need not have generic rank two. The literal global-closure operation also does not give the intended collection of individual Hodge loci.\par
\subsection*{A concrete rank mismatch}
The original four-page Debarre–Laszlo paper [dl] was actually read from the author’s institutional scan. On p. 337 its definition is \(\operatorname{rank}\operatorname{NS}(X)>1\). Remark 2 states that a generic point of a type (b) component has Neron–Severi rank \(g/n\), whereas type (a) has generic rank two. The reviewed Theorem 13 includes type (ii) for every divisor \(n\mid g\), \(n\ne g\).\par
Take \(g=3,n=1\) and a totally real cubic field \(K\). For example, \(x^3+x^2-2x-1\) is irreducible over \(\mathbb Q\) and has positive discriminant \(49\), so its three roots are real. The usual principally polarized Hilbert–Blumenthal construction uses the lattice \(\mathcal O_K\oplus\mathcal D_K^{-1}\) with its perfect alternating trace pairing. In this family \(K\) acts by self-adjoint endomorphisms. The primary source’s identification of rational Neron–Severi classes with self-adjoint endomorphisms gives \[\dim_\mathbb Q\operatorname{NS}(X)_\mathbb Q\ge3\] at every point, and generic rank \(3\). Such a fixed family has no exact-rank-two point. Its appearance in a rank-\(>1\) classification is consistent; its appearance as a component obtained from an exact-rank-two stratum is not.\par
An embedded rank-two lattice can exist inside a rank-three Neron–Severi group. That is a third notion, distinct from both exact rank two and the assertion of rank at least two. It cannot be substituted without changing the definition.\par
\subsection*{Why the entire countable locus cannot be globally closed}
For \(g\ge2\), choose a sufficiently general principally polarized product \(E\times B\), with \(\dim B=g-1\), both factors of Neron–Severi rank one and \(\operatorname{Hom}(E,B)=0\). For \(g=2\), choose two nonisogenous general elliptic curves. The product has Neron–Severi rank exactly two.\par
Its \(\operatorname{Sp}(2g,\mathbb Q)\)-orbit in Siegel space consists of principally polarized abelian varieties isogenous to it: rational symplectic changes give commensurable lattices. Isogeny preserves \(\operatorname{NS}_\mathbb Q\), hence preserves its rank. The rational symplectic group is dense in the real symplectic group, as is seen on a rationally parametrized open Bruhat cell and its translates. The real group acts transitively on Siegel space. Therefore this orbit, and its image in \(\mathcal A_g\), is analytically dense. It is contained in the exact-rank-two locus. Consequently \[\overline{\mathrm{NL}_g^2}^{\,\mathrm{Zar}}=\mathcal A_g
                    \qquad(g\ge2)\] under the literal global-closure definition on p. 668. Its irreducible component is not the collection of proper subvarieties in Theorem 13.\par
A corrected construction should fix rational or integral Hodge conditions, take their closures component by component, and specify the marking and rank convention. It should also identify the descent of the extra Hodge class and the proper morphism used for Chow pushforward. The present check does not claim that the appropriately fixed Hodge-locus construction is impossible.\par
\subsection*{Impact and counterevidence}
This affects Definition 11, the scope of Theorem 13 and the general marked-locus virtual-class discussion. It is a substantive secondary definition/theorem chain, even though the intended construction can plausibly be repaired. The main Torelli/product non-tautology chain does not require this general definition. The product virtual class can be verified directly after the scope is fixed; see below. No failure of that direct product calculation is inferred.\par
\section*{Local corrections with bounded impact}
Problem 3 (Minor: the genus-six displayed coefficient is transposed).  Theorem 5 on p. 665 and the genus-six calculation on p. 716 display \(2370/691\), whereas the definition of \(\Delta_6\) gives \(2730/691\).\par
The exact rational Bernoulli recurrence gives \[B_{12}=-\frac{691}{2730},\qquad
 \frac{6}{6|B_{12}|}=\frac{2730}{691}.\] The coefficients for genera \(1,\ldots,7\) are \(1,10,21,20,11,2730/691,1\). The source’s Hodge integral is \[\int_{\overline{\mathcal M}_6}\lambda_6\lambda_5\lambda_4
                  =\frac{691}{31384184832000}.\] Since the correctly normalized \(\delta_6\) pairs trivially, the incorrectly printed class has pairing with \(\lambda_4\) equal to \[\frac{360}{691}\,
 \frac{691}{31384184832000}
             =\frac{1}{87178291200}\ne0.\] Thus the wrong display is not itself a radical generator. The preserved notebook actually uses \(2730/691\); this is decisive counterevidence limiting the numerical issue to a local correction. It does not resolve Problem Problem 1 (\#prob:certificate). The Bernoulli and rational arithmetic were independently executed in an original reviewer program, not an author notebook.\par
Problem 4 (Minor: sign and reference in the Hodge series).  The definition of the series on p. 677 does not have the sign of the displayed sine expression, and the cited proposition numbers refer to a different primary paper.\par
With the printed convention \(\Lambda(z)=\sum_i z^i\lambda_{g-i}\), one has \[\Lambda(-1)=(-1)^g c(\mathbb E^\vee).\] The genus-one coefficient in the printed series is \(-z/24\), while \[\left(\frac{\sin(t/2)}{t/2}\right)^{-z}\] has coefficient \(+zt^2/24\). Inserting \((-1)^g\) in each genus summand, or using \(c(\mathbb E^\vee)\) instead of \(\Lambda(-1)\), corrects this. The needed generating-series propositions are in the original Hodge integrals and Gromov–Witten theory [fp], not Propositions 3 and 4 of Relative maps and tautological classes, which concern socle and reconstruction. The latter paper remains relevant at its actual uses elsewhere. The corrected series gives the intended positive integral \(J_g\) and does not refute Proposition 2 or Theorem 4.\par
Problem 5 (Major: the leaf patch must fill only its attaching node).  The literal gluing at the end of Section 4.5.6, p. 697, includes too much of the trivial leaf curve when that curve has internal nodes. A restricted open gluing supplies the needed proper Artin family.\par
Let \(T\) be the leaf-side curve, \(q\) its attaching node, and \(r_j\) its internal nodes. The displayed \(T^-\) removes all these nodes. After deleting \(E\) and \(P\), the open \(Z\) has pieces belonging to both \(D_v\) and the other contracted components. Gluing the whole trivial \(T\) along \(T^-\) adds extra smooth representatives over the internal \(r_j\), although the original nodal neighborhoods are already present. The literal special fibre is then nonseparated.\par
Use only the open-and-closed part of \(Z\) belonging to \(D_v\), and glue the trivial family on \(T\setminus\{r_j\}\). The intersection is \(T^-\); only the point \(q\) is filled. The special fibre is \(D_v\), and the resulting flat proper Artin family permits the stated \(\pi_*\mathcal O(Q)=\mathcal O\) argument. This local correction has been carried through in the coefficient-ideal proof. It is not a counterexample to reducedness. The original construction does not produce the proper family required by (39). This step enters Lemma 28 and Theorem 22, and hence the local equations of Proposition 29 used for the excess calculation. Its original proof impact is Major. The verified gluing repair is recorded separately.\par
Problem 6 (Minor: outlook parity and type-(ii) dimensions).  Two assertions on p. 718 need correction to agree with the primary statements actually cited.\par
First, the unqualified claim that \(IH^{2g-2}(\mathcal A_g^{\mathrm{Sat}},\mathbb Q)\) is not generated by \(\lambda\)-classes for all \(g\ge8\) conflicts with Hulek–Tommasi–Taibi, Remark 6(ii) [ht]. That remark states equality with the tautological ring in degree \(2g-2\) for odd \(g\); for odd \(g\ge9\) it instead notes nonzero \(IH^{2g-1}\). Thus \(g=9\) is an explicit discrepancy. Correct the parity/range or give a precise new result superseding the cited one. This concerns an outlook statement; a nonzero Chow class can be homologically trivial, and extra intersection cohomology alone does not identify \(\Delta_g\).\par
Second, the type-(ii) quantity \(g(n+1)/2\) is the dimension, not the codimension. The primary computation is \[\dim S=\frac{g(n+1)}2,\qquad
 \operatorname{codim}_{\mathcal A_g}S=\frac{g(g-n)}2.\] The difference from expected codimension \(g(g-1)/2\) is \(g(n-1)/2\). It is zero for \(n=1\), so the comparison must allow equality. For \(g=2,n=1\), the Hilbert modular surface has dimension \(2\) and codimension \(1\), equal to the expected codimension. These are local formula corrections, distinct from the locus-definition problem above.\par
\subsection*{Other verified local corrections}
The following corrections have bounded, explicit impact; they are not additional Major findings.\par
- On p. 672, when \(g=2k\), the two described fibre-product components remain distinct for the ordered \(\mathcal A_k\times\mathcal A_k\). In \(g=4,k=2\), the distinguished elliptic factor can belong to either labelled factor. Abstract isomorphism does not identify these components. Both excess contributions vanish, so Proposition 17 is unchanged.\par
- On p. 685, an irreducible extremal tree need not have no nontrivial degeneration. A root with two positive-genus leaves can acquire a trivalent genus-zero vertex. The intended statement is that such a tree has no proper smoothing within the extremal trees. The depth definition uses degenerations ending at the given tree and remains consistent.\par
- On p. 687 the immersive map has target \(\mathcal M_g^\mathrm{ct}\) on the chosen local branch, not \(\mathcal A_g\). For \(g=3\), fix an elliptic curve \(E\) and a smooth genus-two curve \(D\), and vary the attaching point on \(D\) away from its Weierstrass points. The stable curve varies while \(E\times J(D)\) is constant. Page 678 already gives the correct target, and the subsequent local chart is a curve deformation chart.\par
- In Example 32 on p. 704 the first numerator must be the total Chern class \(c(N)\), not \(c_1(N)\), to give the displayed excess expression.\par
- On p. 694, the phrase excluding nonnegative powers should exclude positive powers; constants remain. The subsequent coefficient comparison uses positive powers.\par
- On p. 719 the first displayed target of the Hodge map omits the dual. The preceding identification and the following exact sequence correctly use \(\bigwedge^2\mathbb E^\vee\).\par
\section*{Positive verification of the intersection and local-equation chain}
\subsection*{Product loci and the excess bundle}
The product-locus argument uses uniqueness of decomposition of a principally polarized abelian variety into indecomposable principally polarized factors. It does not assume that arbitrary unpolarized isogeny decompositions are unique.\par
For an extremal common refinement of two ordered positive partitions, let \(n_{ij}>0\) denote an occupied row/column cell. Extremality aggregates entries in each cell. The excess bundle is the direct sum of \[\mathbb E_{n_{ij}}^\vee\otimes\mathbb E_{n_{i'j'}}^\vee
       \quad\text{for } i\ne i',\ j\ne j'.\] Its rank is positive when both partitions have at least two parts. Indeed, if an occupied cell \((i,j)\) has no disjoint occupied cell, the existence of another row and another column forces occupied \((i',j)\) and \((i,j')\); these are disjoint from each other.\par
For positive-rank bundles \(V,W\), the symmetric polynomial \(\prod_{a,b}(x_a+y_b)\) lies in the ideal generated by \(c_{\operatorname{rank}V}(V)\) and \(c_{\operatorname{rank}W}(W)\). This follows by setting one root in each family to zero, or from the resultant of the two Chern polynomials when both constant terms vanish. Thus \[e(V)=e(W)=0\quad\Longrightarrow\quad e(V\otimes W)=0.\] The application uses \(\lambda_n=0\) on each \(\mathcal A_n\). It does not require a globally trivial line subbundle. The positive-rank excess bundle consequently has vanishing Euler class, component by component.\par
This is stronger than the vanishing of a class after pushforward. The later boundary restriction of \(\delta_g\) requires precisely this componentwise pullback vanishing. One could not infer \(p^*\alpha=0\) merely from \(p_*p^*\alpha=0\).\par
\subsection*{The only tautological coefficient and the radical}
For \(g\ge2\), the square-free monomial basis of the tautological ring, together with \(\lambda_{g-1}^2=0\), shows that a tautological product class in codimension \(g-1\) annihilated by \(\lambda_{g-1}\) must be a scalar multiple of \(\lambda_{g-1}\).\par
Multiplication of its Torelli pullback by \(\lambda_{g-2}\) eliminates every extremal tree except the single edge between the genus-one root and a leaf of genus \(g-1\). With at least three positive-genus vertices, the largest nonzero Hodge index is at most \(g-3\); with just two such vertices there can be no stable genus-zero vertex in an unmarked tree. The required integrals are \[J_h=\frac{|B_{2h}|}{2h(2h)!},\qquad
 I_g=\frac{|B_{2g}||B_{2g-2}|}
 {2(2g-2)!(2g)(2g-2)}\quad(g\ge2).\] The ratio is exactly \[\frac{J_{g-1}/24}{I_g}
                       =\frac{g}{6|B_{2g}|}.\] For \(g=2\) it is \(10\), consistent with the degree-two ordered product map onto the reduced decomposable divisor.\par
On a separating boundary, both the product-class pullback and the \(\lambda_{g-1}\) pullback vanish: the first by the componentwise excess calculation, the second because every term contains a top Hodge class on one factor. The remaining interior tautological generator has zero pairing by the preceding normalization. This verifies the stated mechanism of Theorem 4, conditional only on the precisely identified standard tautological and Hodge-integral results. It does not supply nonvanishing.\par
\subsection*{The scheme structure, not only its reduced support}
The local analysis fixes a domain branch corresponding to a chosen elliptic factor; when several factors exist, they must not be merged silently. The degree-one stable-map description and the Jacobian product description agree in families under the checked Cartesian statement of Greer–Lian [gl]. The degree-one map is an isomorphism on the elliptic root and contracts the complementary compact-type curve in the closed fibre.\par
In a boundary-normal calculation, evaluation of canonical forms at a smooth attaching point on a positive-genus component is nonzero by basepoint freeness. This also holds for hyperelliptic components; no assertion that the global Torelli map is immersive there is needed.\par
For an extremal tree, let \(a_e\) be the independent node-smoothing parameters. The reduced support has the square-free path ideal \[I=\left(\prod_{e\in P_v}a_e:
                   v\text{ is a positive-genus leaf}\right),\] where \(P_v\) is the path from the distinguished root to \(v\). Every such path must retain a node, or the root would merge with positive genus. Square-freeness establishes reducedness of this candidate monomial scheme, but alone does not establish that it is the actual scheme.\par
The following Artin check supplies that additional step. Work over a local Artin complex algebra \((A,m)\), with \(m^N=0\), and use compatible plumbing charts \[xy=a,\qquad a\in m.\] A regular function has normal form \[c+\sum_{i\ge1}p_i x^i+\sum_{j\ge1}q_j y^j.\] On the punctured \(x\)-branch its \(x^{-j}\) coefficient is \(a^j q_j\); \(q_j\) is the actual positive Taylor coefficient on the opposite branch. This identity does not divide by a nilpotent. Expressions such as \(s_j/a^j\) in the paper must mean these compatible coefficients.\par
For a proper compact-type leaf \((D,q)\), with \(q\) on a positive-genus component, \[H^0(D,\mathcal O_D(q))=\mathbb C.\] A function with at most one simple pole on that component would otherwise give a degree-one map to \(\mathbb P^1\); the other branches contribute no sections vanishing at their attachment. The same assertion over a flat proper Artin family is \(\pi_*\mathcal O(Q)=A\): induction through a small extension reduces any section to a constant and then uses \(J\otimes H^0(\mathcal O_D(q))=J\). The restricted completion in Problem Problem 5 (\#prob:patch) supplies the needed proper family.\par
Propagate the root coordinate over the rational skeleton, obtaining \(G\), and subtract it from the actual pulled-back target coordinate \(F\). Let \(J\subset A\) be generated by all principal-part coefficients of \(F-G\) on internal rational components and by all root-to-leaf products. At an internal edge the node identity puts every principal-part coefficient in \(mJ\). At a leaf, modulo \(mJ\), \(F\) has only a possible simple pole with unit times its path product as coefficient; all higher poles have disappeared. The proper-leaf fact forces this pole and all nonconstant Taylor coefficients to vanish. The reverse node identity then puts the remaining parent-side coefficients in \(mJ\). Therefore \[J\subseteq mJ,
            \qquad J\subseteq m^N J=0.\] This simultaneous argument covers arbitrary finite trees and direct root-to-leaf edges. It allows the elliptic target to deform; it does not require the unexpanded global analytic product assertion on p. 688.\par
If \(K\) is the actual defining ideal, the reduced-support calculation gives \(K\subseteq I\). The Artin argument shows that every generator of \(I\) vanishes in every Artin quotient of the completed local ring modulo \(K\). Noetherian Krull intersection gives \(I\subseteq K\), hence \(I=K\). The complete supporting derivation, including coordinate and properness assumptions, is retained in the verification evidence.\par
An original finite combinatorial check examined 66 admissible trees with two through nine vertices, 10,278 edge-support patterns and 184 minimal cutsets. It verified the path-ideal and cutset identities in those cases. The all-tree argument above, not these finite examples, is the basis for the general conclusion.\par
\subsection*{Excess recursion and sample calculations}
In Section 5, put \(n=|E(T)|\), let \(k\) count the positive-genus leaves, and put \(r=g-1\). Stability and positive leaf genera give \(n\ge k\) and \(r\ge k\). A contribution has degree \(r-n\); hence \(n>r\) contributes zero. For a star the normal-bundle quotient gives \[\left[\frac{c(N)}{\prod_e(1+z_e)}\right]_{r-k},
             \qquad z_e=-\psi_e'-\psi_e''.\] The elliptic-root Hodge line has zero first Chern class on \(\mathcal A_1\). On the equivariant coordinate model, subtracting the previously determined smoothing-stratum contributions leaves a class supported at the origin. Its pushforward is divisible by \(\prod_e z_e\). Multiplication by this product is injective in the polynomial ring of independent equivariant parameters, justifying division in (41). The degree \(r-n\le r-k\) also ensures unique expression through the required Chern classes of \(N\). The underlying trivial lines on p. 699 carry the path characters specified by the displayed equivariant Chern formula; they are not weight-zero equivariant lines.\par
The global application uses the established monomial local model and residual-intersection functoriality, not agreement of normal bundles alone. The actual Fulton statements were checked on pp. 161–164 of the second edition [fulton]: Theorem 9.2 uses a regular embedding and a Cartier divisorial part; Corollary 9.2.1 requires a regular residual embedding; Corollary 9.2.3 uses blow-up and the residual Segre class. It does not declare an arbitrary residual scheme regular. The coordinate-stratum application respects these distinctions.\par
An original reviewer symbolic calculation reproduced Examples 30 and 31, the corrected \(-5\) in Example 32, the expressions in Example 34, and \(15\) in Example 35. It enumerated the edge-subset contractions for Examples 30–35, tested Remark 33 for \(k=2,\ldots,6\), and checked the coefficient ratio for \(g=2,\ldots,10\). There are 26 exact matches and the two deliberate discrepancies: the literal \(c_1(N)\) numerator and the genus-one series sign. These are independent small symbolic checks, not admcycles execution.\par
Graph automorphisms fix the root and preserve genera. The factor \(1/|\operatorname{Aut}(T)|\) in (44) accounts for the labelled graph chart. In genus two the rooted one-edge graph has factor \(1\), while its gluing map has degree \(2\) onto the reduced separating divisor. The genus-six assembly contains all 24 displayed contributions; principal automorphism factors \(1,1,1,2,2,6,120\) and triple-intersection multiplicities agree with the static notebook after vertex relabelling. An extra \(-\lambda_3\) on a compact-type factor vanishes there. These checks establish transcription and local identities, not the full quotient reduction needed in Problem Problem 1 (\#prob:certificate).\par
\section*{Genus seven, cohomology and virtual classes}
\subsection*{Genus seven, conditional on the nonzero input}
Let \[\alpha=j_*(\pi^*\delta_6\times[\mathcal M_{1,1}])\in R^6(\mathcal M_7^\mathrm{ct}).\] Elliptic-tail gluing is not globally an embedded divisor: additional branches occur when the genus-six factor already has an elliptic tail. Here those terms vanish because the boundary pullbacks of \(\delta_6\) are zero, by Theorem 4’s proof. The remaining self-intersection term is \(-\psi_1\pi^*\delta_6\), consistent with the extra-boundary term in CLS formula (19). The forgetful pushforward gives \[\pi_*(-\psi_1\pi^*\delta_6)=-10\delta_6.\] External product with \([\mathcal M_{1,1}]\) is injective: pullback to a chosen elliptic-curve point is a left inverse with rational Chow coefficients. The socle/gluing calculation places \(\alpha\) in the one-dimensional genus-seven radical. Conditional on \(\delta_6\ne0\), its boundary pullback is nonzero. The boundary pullback of \(\delta_7\) is zero, so \(\delta_7=c\alpha\) forces \(c=0\). There is no separate self-intersection defect once this hypothesis is explicit. The deduction still depends on the genus-six nonzero step. In lower genera, a zero Torelli pullback likewise does not establish equality of the original classes on \(\mathcal A_g\).\par
\subsection*{The cohomological range}
The checked Hulek–Tommasi version with Taibi’s appendix [ht] gives \(IH^*(\mathcal A_g^{\mathrm{Sat}},\mathbb Q)=R_g\) for \(g\le5\). The actual statements identify the tautological classes, not just an abstract dimension match. Theorem 33 gives the analogous description for \(k<2g-2\), at level one; no arbitrary-level claim is inferred. Durfee’s Lemma 2(2) [durfee], read in the primary 1993 author version, supplies the weight-graded surjection \[\operatorname{Gr}^W_kIH^k(X)\longrightarrow
                       \operatorname{Gr}^W_kIH^k(U)\] for an open inclusion \(U\subset X\). Take \(X=\mathcal A_g^{\mathrm{Sat}}\) and \(U=\mathcal A_g\). On the smooth open stack, rational intersection cohomology agrees with ordinary rational cohomology; finite stabilizers introduce no additional rational contribution. Codimension-\(r\) algebraic classes have weight \(2r\). Thus Proposition 7’s cohomological tautologicity follows for \(g\le5\), and in codimension \(r<g-1\) generally. This is not equality in Chow. Problem Problem 6 (\#prob:outlook) concerns the later endpoint assertion.\par
\subsection*{The direct product virtual class}
Let \(j:\mathcal A_k\times\mathcal A_{g-k}\to\mathcal A_g\). Writing \(p_1,p_2\) for the projections, the formulas are \[\begin{aligned}
 N_j&=p_1^*\mathbb E_k^\vee\otimes p_2^*\mathbb E_{g-k}^\vee,\\
 \operatorname{Obs}
   &=p_1^*(\bigwedge^2\mathbb E_k^\vee)\oplus
                              p_2^*(\bigwedge^2\mathbb E_{g-k}^\vee).\end{aligned}\] The ranks satisfy \(k(g-k)+\binom{k}{2}+\binom{g-k}{2}=\binom g2\). The staircase Schur-determinant induction in Lemma 41 uses \(\lambda_g=0\) and \(\lambda_{g-1}^2=0\) at the indicated steps. Dualization gives the signs \((-1)^{\binom{k}{2}}\) and \((-1)^{\binom{g-k}{2}}\) in Proposition 14. This direct calculation survives the general-locus definition correction. Speculation 15 remains a speculation.\par
\section*{Dependency boundaries and disposition}
Core references were read at their actually used statements:\par
- Clemens–Griffiths, original 1972 publisher scan, Definition 3.5, Lemma 3.20 and Corollary 3.23 (pp. 293, 296–297), supports the polarized indecomposable factorization in Proposition 16 [cg].\par
- Van der Geer’s 1996 v2 author version, not an asserted byte-identical 1999 publication, gives the Hodge relation, \(\lambda_g=0\), and square-free basis in Theorem 1.1, Proposition 1.2 and Theorem 1.5 [vdg].\par
- The actual Faber–Pandharipande author manuscript gives the Hodge integrals [fp]; the separately read publisher Relative maps paper supplies the compact-type/rational-tail socle input [relative]. The nonzero triple-Hodge integral identifies \(\lambda_{g-2}\) as the interior top-degree generator. Ordinary Chow localization handles the boundary-supported remainder; no conjectural exact tautological restriction sequence is used.\par
- Greer–Lian’s v3 family statement and Sernesi’s Artin rigidity theorem are applied under their local-family hypotheses [gl, sernesi]. The original coefficient-ideal derivation supplies the general-tree bookkeeping.\par
- The actual CLS PDF supports the quotient and radical dimensions, but its reference back to this paper does not certify the specified nonzero vector [cls]. The Debarre–Laszlo, Hulek–Tommasi–Taibi and Durfee primary versions support the scope comparisons [dl, ht, durfee].\par
This is not a reproof of all external intersection theory, mixed-Hodge theory, automorphic classification or relation-rank computations. Exact versions, read ranges, image inspections, retrieval records and hashes are retained. A checked theorem application is distinguished from an independently reproduced proof.\par
There are two Major findings, three numbered Minor findings, one Needs Evidence finding, and six bounded local clarifications. The substantive nonzero certificate remains to be supplied; the exact-rank/global-closure definition needs correction. The original leaf gluing is invalid in a core proof, and its verified repair is recorded separately. The numerical, sign, and outlook corrections have the limited effects described above.\par
The preserved original and recovery records attribute the full-source reading, original arithmetic, primary checks and bounded helper work to their respective reviewers. The present independent focused acceptance rechecks the reported findings and limits; its actual current TeX build and final-page inspections are separately bound in the verification files. Historical checks retain their original attribution and are not relabelled as this acceptance. No author notebook was executed or software installed. The retained notebook outputs remain author-supplied observations. Scientific decision is pending; human reviews are empty. The original and recovery tokens are, respectively,\par
99 S. Canning, D. Oprea and R. Pandharipande, Tautological and non-tautological cycles on the moduli space of Abelian varieties, Invent. math. 242 (2025), 659–723. Assigned publisher PDF. https://doi.org/10.1007/s00222-025-01367-4. S. Canning, H. Larson and J. Schmitt, The Gorenstein property and Pixton’s conjecture for compact type moduli. Cited 29-page author PDF, https://people.math.ethz.ch/\textasciitilde{}scanning/compacttype.pdf. Observed file metadata: October 2024; historical byte identity is not asserted. The separately checked https://arxiv.org/abs/2607.02249v2 is a later 2026 version. O. Debarre and Y. Laszlo, Le lieu de Noether–Lefschetz pour les varietes abeliennes, C. R. Acad. Sci. Paris Ser. I Math. 311 (1990), 337–340. Actual institutional scan: https://perso.imj-prg.fr/wp-content/uploads/debarre-pub/19.pdf. C. Faber and R. Pandharipande, Hodge integrals and Gromov–Witten theory. Actual 25-page author manuscript dated 1 June 1999: https://people.math.ethz.ch/\textasciitilde{}rahul/hodge.pdf. C. Faber and R. Pandharipande, Relative maps and tautological classes. Actual 2005 publisher PDF: https://ems.press/content/serial-article-files/31534. K. Hulek and O. Tommasi, with appendix by O. Taibi, The topology of \(\mathcal A_g\) and its compactifications. Actual arXiv:1712.03771v3, 15 May 2018; Theorems 17, 32, 33 and Remark 6(ii). https://arxiv.org/abs/1712.03771v3. A. H. Durfee, Intersection homology Betti numbers. Actual arXiv:alg-geom/9302004v1, 12 February 1993, Lemma 2; not a claimed byte-identical 1995 journal PDF. https://arxiv.org/abs/alg-geom/9302004v1. W. Fulton, Intersection Theory, second edition, Springer, 1998. Actual scanned pp. 161–164, Theorem 9.2 and Corollaries 9.2.1, 9.2.3. Only these publicly accessible page images were obtained; the whole book was not claimed read. C. H. Clemens and P. A. Griffiths, The intermediate Jacobian of the cubic threefold, Ann. of Math. 95 (1972), 281–356. Actual publisher scan from the IAS publication site; Definition 3.5, Lemma 3.20 and Corollary 3.23 checked. G. van der Geer, Cycles on the moduli space of abelian varieties. Actual arXiv:alg-geom/9605011v2, 17 June 1996. The later 1999 publication’s typesetting and page numbering were not compared. https://arxiv.org/abs/alg-geom/9605011v2. F. Greer and C. Lian, \(d\)-elliptic loci and the Torelli map. Actual arXiv:2404.10826v3, 8 September 2024, Proposition 3.12 and its Cartesian family proof. https://arxiv.org/abs/2404.10826v3. E. Sernesi, Deformations of Algebraic Schemes, Springer, 2006. Actual chapter-one publisher sample, Theorem 1.2.4 on smooth affine rigidity over Artin bases; not a claim to have read the whole book.\par
\end{document}
